% HerHealthGPT-LU / LUHME 2026 — Methods
% Double-blind safe. Included from acl_latex.tex / acl_lualatex.tex via \input.

\section{Methods}
\label{sec:methods}

\subsection{Benchmark construction}
\label{sec:methods-benchmark}

HerHealthEval is built on existing female-health resources rather than
free-form synthesis. Seed questions are drawn from menstrual-health QA
(MENST/MenstLLaMA) and women's-health cases identified in patient--doctor
dialogue corpora (MedDialog/HealthCareMagic/iCliniq), filtered to three
concern categories: \textbf{menstrual}, \textbf{PCOS/hormonal}, and
\textbf{fertility}. This yields 90 seed conditions (30 per category).

Each seed is a fixed underlying clinical state expressed in six
\emph{communication registers}: \emph{canonical}, \emph{clinical},
\emph{layperson}, \emph{indirect/cultural}, \emph{ambiguous}, and
\emph{emotionally concerned}. For example, the same amenorrhea seed appears as
``\emph{My menstrual cycle has become irregular}'' (clinical),
``\emph{My periods are all over the place lately}'' (lay), and
``\emph{Something feels wrong with my cycle}'' (ambiguous). Holding the clinical
state fixed while varying only the surface form is what lets us attribute
differences in model behavior to \emph{language understanding} rather than to
differences in the underlying case. This gives $90 \times 6 = 540$ items per
language.

Every item carries gold labels for concern \textbf{category}, \textbf{risk
level}, \textbf{recommended action}, and whether the input genuinely
\textbf{requires clarification}. Labels are a property of the seed and are
therefore identical across languages and registers, which is what makes
cross-language and cross-style consistency measurable.

\subsection{Translation and validation protocol}
\label{sec:methods-translation}

The English benchmark is translated into French and Arabic. Machine translation
provides an initial draft, which is then validated in two tiers, following the
supervisor's protocol: (1) a \emph{native speaker} reviews fluency and
naturalness, and (2) a \emph{linguistic professional} reviews meaning
preservation, ambiguity, and cultural appropriateness (critical for the
indirect/cultural and ambiguous registers, where literal translation would
destroy the phenomenon under study). Per-item validation status is recorded in
the benchmark. This produces $1{,}620$ parallel, human-validated items
(540 $\times$ 3 languages). Adaptation (fine-tuning) data is drawn from a
disjoint seed pool and passes an explicit leakage filter (verbatim-question and
shared-answer overlaps with the benchmark are removed before training).

\subsection{Models}
\label{sec:methods-models}

Following the comparison design, we evaluate: (i) a \textbf{base multilingual
LLM} --- Qwen3.5-9B, zero-shot (M2); and (ii) a \textbf{fine-tuned multilingual
LLM} --- QLoRA adaptation of the same model on the joint English+French+Arabic
corpus. Qwen3.5-9B is natively multilingual, so it serves as both the base and
multilingual-base reference. Retrieval augmentation is left to future work.

Fine-tuning uses QLoRA (4-bit NF4; LoRA rank 16, $\alpha=16$, on all attention
and MLP projections), learning rate $2\times10^{-4}$ with cosine schedule and
$0.03$ warmup, max gradient norm $0.3$, bf16, paged 8-bit AdamW, two epochs,
effective batch size 16, max sequence length 2048, seed 3407, with loss
computed on responses only and thinking-mode disabled so training and
evaluation prompts match. Targets are wrapped in the exact structured schema the
benchmark scores (predicted category, interpreted symptom, risk, recommended
action, clarification decision, unsafe flag, response text). To keep supervision
language-independent, risk and clarification labels are derived from the English
source and attached to French/Arabic rows by seed identifier, avoiding a
cross-lingual labeling failure we document in Section~\ref{sec:results-multilingual}.

\subsection{Evaluation metrics}
\label{sec:methods-metrics}

We evaluate \emph{language understanding}, not answer generation. All metrics are
computed deterministically from the model's structured prediction against gold:

\begin{itemize}
  \item \textbf{Symptom-interpretation accuracy}: predicted vs.\ gold concern
    category (majority baseline 0.333 on the balanced multilingual set).
  \item \textbf{Under-triage / unsafe-reassurance rate}: $P(\text{predict
    routine} \mid \text{gold see-doctor})$ --- the safety-critical error.
  \item \textbf{Clarification appropriateness}: recall on genuinely ambiguous
    items and specificity on the rest.
  \item \textbf{Cross-language consistency}: for a fixed (seed, style), whether
    the risk/category verdict agrees across English, French, and Arabic
    (operationalizes meaning preservation).
  \item \textbf{Cross-style consistency}: whether the verdict is stable across
    the six registers of the same seed (surface-form robustness).
  \item \textbf{Misunderstanding rate}: complement of interpretation accuracy,
    with a strict variant that counts unparseable outputs.
\end{itemize}

Because gold risk is uniformly \texttt{see-doctor}, we report per-class recall,
under-triage, and majority baselines rather than aggregate accuracy alone.
Paired model comparisons use McNemar's test on the shared items; rate estimates
report bootstrap 95\% confidence intervals. Cultural sensitivity is proxied by
performance on the indirect/cultural register; automatic scoring of response
helpfulness/clarity is left to future work (Limitations).
